\documentclass[12pt,a4paper]{article}

\usepackage[utf8]{inputenc}
\usepackage[spanish]{babel}
\usepackage[a4paper,margin=2.5cm]{geometry}
\usepackage{longtable}
\usepackage{array}
\usepackage{hyperref}
\usepackage{verbatim}

\hypersetup{
    colorlinks=true,
    linkcolor=black,
    urlcolor=blue,
    pdftitle={Informe de simulacion MIL-STD-1553 sobre SPI},
    pdfauthor={Proyecto Raspberry Pi Pico 2 W}
}

\title{Informe Tecnico\\Simulacion Logica de MIL-STD-1553 sobre SPI}
\author{Proyecto basado en Raspberry Pi Pico 2 W}
\date{\today}

\begin{document}

\maketitle

\tableofcontents
\newpage

\section{Introduccion}

El presente documento consolida y deja documentado el estado actual del desarrollo de una
simulacion logica de MIL-STD-1553 sobre SPI utilizando dos placas Raspberry Pi Pico 2 W y un
dashboard provisional en Flask. El objetivo general de esta etapa fue evolucionar desde un esquema
de mini-tramas SPI simples hacia una representacion mucho mas fiel, a nivel logico, del
comportamiento del estandar MIL-STD-1553, manteniendo todavia una implementacion de bajo
riesgo, de facil depuracion y sin entrar aun en la capa fisica real del bus.

Durante esta fase se consolidaron tres piezas principales:

\begin{itemize}
    \item un emisor \textbf{MIL-M} que construye palabras y mensajes logicos 1553;
    \item un receptor \textbf{MIL-E} que valida, interpreta contexto de mensaje y emite eventos
    estructurados por serial;
    \item un \textbf{dashboard Flask} provisional que consume la salida del esclavo y la presenta
    como sniffer de prueba mientras no se dispone de una Raspberry dedicada como host final.
\end{itemize}

La simulacion lograda no es aun un bus MIL-STD-1553 fisico real. No hay codificacion Manchester,
ni sincronismos fisicos de palabra, ni front-end analogico 1553, ni tiempos de linea de 1 Mb/s
medidos sobre el medio real. Sin embargo, a nivel conceptual y de flujo de protocolo, el sistema
ya representa de manera util:

\begin{itemize}
    \item palabras \texttt{COMMAND}, \texttt{DATA} y \texttt{STATUS};
    \item mensajes \texttt{BC->RT}, \texttt{RT->BC}, \texttt{BROADCAST} y \texttt{MODE CODE};
    \item chequeos de integridad;
    \item deteccion de anomalias de secuencia y de orden de mensaje;
    \item una salida estructurada apta para monitoreo y analisis en host.
\end{itemize}

\section{Objetivos de la etapa}

Los objetivos concretos perseguidos en esta etapa fueron los siguientes:

\begin{itemize}
    \item reemplazar el esquema inicial de mini-tramas manuales por una representacion logica
    inspirada en MIL-STD-1553;
    \item mantener SPI como transporte interno entre placas, priorizando robustez y facilidad de
    prueba;
    \item separar responsabilidades entre emisor, receptor y host;
    \item aumentar el realismo del protocolo antes de migrar a la capa fisica real;
    \item dejar trazabilidad de errores, warnings y secuencias de mensaje;
    \item disponer de una herramienta de visualizacion temporal en PC mediante Flask.
\end{itemize}

\section{Alcance y limites}

\subsection{Que se simula}

Actualmente se simulan de manera logica los siguientes conceptos del estandar:

\begin{itemize}
    \item estructura de palabra de comando;
    \item estructura de palabra de estado;
    \item palabra de datos de 16 bits;
    \item bit de direccion \texttt{T/R};
    \item \texttt{RT Address};
    \item \texttt{Subaddress};
    \item \texttt{Word Count};
    \item \texttt{Mode Codes};
    \item direccion de broadcast;
    \item secuencias \texttt{BC->RT} y \texttt{RT->BC};
    \item validacion de orden esperado en el receptor;
    \item estados simples del terminal remoto.
\end{itemize}

\subsection{Que todavia no se simula fisicamente}

Todavia no se implementaron los siguientes aspectos de una solucion MIL-STD-1553 real:

\begin{itemize}
    \item codificacion Manchester II;
    \item sincronismos fisicos de palabra;
    \item temporizacion real de 20 us por palabra;
    \item ventana real de respuesta RT entre 4 us y 12 us;
    \item acople por transformador;
    \item front-end analogico diferencial;
    \item buses redundantes fisicos A/B;
    \item captura pasiva sobre una linea 1553 real.
\end{itemize}

\section{Arquitectura resultante}

La arquitectura quedo dividida en tres componentes claramente separados:

\begin{longtable}{|>{\raggedright\arraybackslash}p{3cm}|>{\raggedright\arraybackslash}p{10.8cm}|}
\hline
\textbf{Bloque} & \textbf{Responsabilidad actual} \\
\hline
MIL-M & Construccion y transmision de palabras y mensajes logicos MIL-STD-1553 sobre SPI. \\
\hline
MIL-E & Recepcion byte a byte, validacion de integridad, seguimiento de contexto de mensaje y emision de eventos estructurados por serial. \\
\hline
Flask-mil & Dashboard provisional en PC para consumir la salida serial del esclavo y actuar como sniffer temporal. \\
\hline
\end{longtable}

Esta separacion fue una decision deliberada de arquitectura. Se evito convertir al esclavo en un
dashboard completo y tambien se evito sobrecargar al master con responsabilidades de analisis.
El resultado es una base mas modular y escalable.

\section{Formato de mini-trama SPI}

Cada palabra logica MIL-STD-1553 se encapsula actualmente en una mini-trama SPI fija de 7 bytes:

\begin{center}
\texttt{[SYNC][TYPE][WORD\_MSB][WORD\_LSB][PARITY][SEQ][CHECKSUM]}
\end{center}

\subsection{Descripcion de campos}

\begin{longtable}{|>{\raggedright\arraybackslash}p{2.7cm}|>{\raggedright\arraybackslash}p{1.8cm}|>{\raggedright\arraybackslash}p{8.7cm}|}
\hline
\textbf{Campo} & \textbf{Tamano} & \textbf{Descripcion} \\
\hline
\texttt{SYNC} & 1 byte & Byte fijo de sincronizacion. Se utiliza \texttt{0xAA}. \\
\hline
\texttt{TYPE} & 1 byte & Tipo de palabra: \texttt{0x01} para \texttt{COMMAND}, \texttt{0x02} para \texttt{DATA}, \texttt{0x03} para \texttt{STATUS}. \\
\hline
\texttt{WORD\_MSB} & 1 byte & Byte mas significativo de la palabra logica de 16 bits. \\
\hline
\texttt{WORD\_LSB} & 1 byte & Byte menos significativo de la palabra logica de 16 bits. \\
\hline
\texttt{PARITY} & 1 byte & Paridad impar calculada sobre los 16 bits de la palabra. \\
\hline
\texttt{SEQ} & 1 byte & Contador secuencial por palabra transmitida. \\
\hline
\texttt{CHECKSUM} & 1 byte & XOR de los 6 bytes anteriores como chequeo rapido de transporte. \\
\hline
\end{longtable}

\subsection{Motivacion del formato}

El formato se mantuvo deliberadamente pequeno y estable para:

\begin{itemize}
    \item conservar la robustez del envio byte a byte;
    \item simplificar la recepcion en el esclavo;
    \item facilitar la depuracion por consola;
    \item permitir evolucionar la semantica del protocolo sin cambiar aun el transporte.
\end{itemize}

\section{Modulo compartido de protocolo}

Se creo y consolidó un modulo comun de protocolo:

\begin{itemize}
    \item \texttt{mil1553\_protocol.h}
    \item \texttt{mil1553\_protocol.c}
\end{itemize}

Este modulo centraliza la logica compartida por master y slave, evitando duplicacion y asegurando
consistencia en el tratamiento del protocolo.

\subsection{Funciones incluidas}

\begin{itemize}
    \item calculo de checksum de transporte;
    \item calculo de paridad impar;
    \item construccion de \texttt{COMMAND WORD};
    \item construccion de \texttt{DATA WORD};
    \item construccion de \texttt{STATUS WORD};
    \item encapsulado de mini-trama SPI;
    \item validacion de mini-trama recibida;
    \item decodificacion de \texttt{COMMAND WORD};
    \item decodificacion de \texttt{STATUS WORD};
    \item deteccion de \texttt{Mode Code};
    \item deteccion de \texttt{Broadcast};
    \item calculo de \texttt{Word Count} efectivo, incluyendo la regla \texttt{0 => 32}.
\end{itemize}

\subsection{Modelo de Command Word}

La palabra de comando implementada sigue la estructura logica:

\begin{itemize}
    \item \texttt{RT Address}: 5 bits;
    \item \texttt{T/R}: 1 bit;
    \item \texttt{Subaddress}: 5 bits;
    \item \texttt{Word Count} o \texttt{Mode Code}: 5 bits.
\end{itemize}

\subsection{Modelo de Status Word}

La palabra de estado implementada incluye los siguientes campos:

\begin{itemize}
    \item \texttt{RT Address};
    \item \texttt{message error};
    \item \texttt{instrumentation};
    \item \texttt{service request};
    \item \texttt{broadcast received};
    \item \texttt{busy};
    \item \texttt{subsystem flag};
    \item \texttt{dynamic bus acceptance};
    \item \texttt{terminal flag}.
\end{itemize}

\section{Evolucion funcional de MIL-M}

\subsection{Primera etapa del emisor}

El master comenzo emitiendo secuencias simples del estilo:

\begin{center}
\texttt{COMMAND + DATA + STATUS}
\end{center}

en forma repetitiva y con carga util incremental.

\subsection{Estado actual del emisor}

En su estado actual, \textbf{MIL-M} ya puede generar distintos tipos de mensaje logico:

\begin{itemize}
    \item \texttt{BC->RT};
    \item \texttt{RT->BC};
    \item \texttt{BROADCAST};
    \item \texttt{MODE\_CODE}.
\end{itemize}

\subsection{Mensajes BC a RT}

La secuencia implementada es:

\begin{center}
\texttt{COMMAND + N DATA + STATUS}
\end{center}

Esto simula un controlador de bus ordenando a un RT recibir una cantidad variable de palabras.

\subsection{Mensajes RT a BC}

La secuencia implementada es:

\begin{center}
\texttt{COMMAND + STATUS + N DATA}
\end{center}

Esto representa a un RT respondiendo con estado y luego con carga util de datos.

\subsection{Mensajes Broadcast}

Se incorporo soporte para mensajes con \texttt{RT Address = 31}, tratados como broadcast. En
esta simulacion se utiliza:

\begin{center}
\texttt{COMMAND + N DATA}
\end{center}

sin palabra de estado posterior, respetando la logica general de evitar respuestas en broadcast.

\subsection{Mode Codes}

Tambien se incorporaron mensajes de \texttt{Mode Code}, usando \texttt{Subaddress = 0} o
\texttt{31}. En la simulacion actual se emplea una variante simple del tipo:

\begin{center}
\texttt{COMMAND + STATUS}
\end{center}

Esto permite comenzar a ejercitar el camino logico de control y supervision del RT.

\subsection{Word Count variable}

Se agrego soporte para \texttt{word count} variable por canal. El master puede transmitir mensajes
con varias palabras de datos, en vez de limitarse a una sola.

\subsection{Canales logicos de prueba}

Se utilizaron canales de simulacion para hacer las pruebas mas legibles:

\begin{itemize}
    \item temperatura;
    \item velocidad;
    \item presion;
    \item presion broadcast;
    \item mode sync.
\end{itemize}

\subsection{Logging de MIL-M}

El emisor imprime trazas de depuracion por cada palabra transmitida, indicando:

\begin{itemize}
    \item secuencia \texttt{SEQ};
    \item tipo de palabra;
    \item palabra cruda;
    \item paridad;
    \item frame SPI completo;
    \item interpretacion basica de la palabra.
\end{itemize}

Esto fue clave para validar visualmente la construccion correcta del protocolo.

\section{Evolucion funcional de MIL-E}

\subsection{Primera etapa del receptor}

El esclavo comenzo como receptor minimo de transporte. Su funcion era:

\begin{itemize}
    \item reconstruir la mini-trama SPI;
    \item validar \texttt{SYNC};
    \item validar paridad;
    \item validar \texttt{CHECKSUM};
    \item imprimir la palabra en una salida serial legible.
\end{itemize}

\subsection{Estado actual del receptor}

En su estado actual, \textbf{MIL-E} hace bastante mas que validacion basica. Manteniendose dentro
de su rol de receptor/transporte, ahora implementa:

\begin{itemize}
    \item validacion de integridad por palabra;
    \item deteccion de saltos de secuencia \texttt{SEQ};
    \item seguimiento de contexto de mensaje;
    \item deteccion de \texttt{Mode Code};
    \item deteccion de \texttt{Broadcast};
    \item conteo esperado y observado de palabras de datos;
    \item deteccion de anomalias de orden en la secuencia de palabras;
    \item emision de eventos estructurados para el host.
\end{itemize}

\subsection{Eventos seriales emitidos}

La interfaz serial del esclavo fue estandarizada en un conjunto de eventos aptos para ser
consumidos por un host o dashboard:

\begin{longtable}{|>{\raggedright\arraybackslash}p{3.1cm}|>{\raggedright\arraybackslash}p{8.9cm}|}
\hline
\textbf{Evento} & \textbf{Uso} \\
\hline
\texttt{RX\_OK} & Frame valido recibido y parseado. \\
\hline
\texttt{RX\_ERR} & Frame invalido con detalle de errores de integridad. \\
\hline
\texttt{RX\_WARN} & Warning de transporte o de protocolo. \\
\hline
\texttt{RX\_MSG\_START} & Inicio de un nuevo mensaje logico. \\
\hline
\texttt{RX\_DATA} & Observacion de una palabra de datos dentro del mensaje activo. \\
\hline
\texttt{RX\_STATUS} & Observacion de una palabra de estado dentro del mensaje activo. \\
\hline
\texttt{RX\_MSG\_END} & Fin del mensaje logico, con resultado \texttt{OK} o \texttt{WARN}. \\
\hline
\end{longtable}

\subsection{Warnings logicos implementados}

Ademas de la validacion de transporte, el esclavo puede detectar y reportar:

\begin{itemize}
    \item \texttt{SEQ\_GAP};
    \item \texttt{COMMAND\_BEFORE\_PREVIOUS\_END};
    \item \texttt{STATUS\_IN\_BROADCAST};
    \item \texttt{STATUS\_BEFORE\_ALL\_DATA};
    \item \texttt{DUPLICATE\_STATUS};
    \item \texttt{DATA\_AFTER\_STATUS};
    \item \texttt{DATA\_OVERFLOW};
    \item \texttt{ORPHAN\_STATUS};
    \item \texttt{ORPHAN\_DATA};
    \item \texttt{UNKNOWN\_TYPE}.
\end{itemize}

\subsection{Flujo de arranque del esclavo}

Se agregaron mensajes \texttt{[INFO]} al iniciar, y tambien se incorporo \texttt{fflush(stdout)}
tras cada evento importante para que la salida llegue de forma mas estable al lector serial del host.

\section{Dashboard Flask provisional}

\subsection{Motivacion}

Como todavia no se dispone de una Raspberry dedicada para actuar como host final, se decidio
adaptar la base existente en Flask que originalmente estaba orientada a ARINC, transformandola en
un dashboard temporal para MIL.

\subsection{Objetivo del dashboard}

Su objetivo es actuar como sniffer provisional, permitiendo verificar rapidamente que:

\begin{itemize}
    \item el esclavo recibe correctamente;
    \item los frames validos e invalidos se contabilizan;
    \item los mensajes logicos se reconstruyen;
    \item los warnings se muestran;
    \item la salida serial del esclavo es apta para analisis en host.
\end{itemize}

\subsection{Cambios realizados en Flask}

El backend \texttt{app.py} fue completamente adaptado para:

\begin{itemize}
    \item leer por puerto serial;
    \item parsear eventos \texttt{RX\_...};
    \item mantener buffers de consola, frames, warnings y mensajes;
    \item reconstruir mensajes a partir de \texttt{RX\_MSG\_START}, \texttt{RX\_DATA},
    \texttt{RX\_STATUS} y \texttt{RX\_MSG\_END};
    \item exponer estadisticas y endpoints JSON para la interfaz web.
\end{itemize}

\subsection{Visualizaciones implementadas en la interfaz}

La interfaz HTML fue redise\~nada para mostrar:

\begin{itemize}
    \item consola serial;
    \item frames validos e invalidos;
    \item warnings recientes;
    \item mensajes reconstruidos;
    \item resumen por tipo de palabra;
    \item resumen por direccion del mensaje;
    \item resumen por resultado;
    \item ultima palabra de estado observada.
\end{itemize}

\subsection{Uso del filtro}

El dashboard incluye un filtro textual simple que actua sobre consola y log. Por ejemplo, permite
filtrar lineas que contengan:

\begin{itemize}
    \item \texttt{RX\_OK};
    \item \texttt{RX\_WARN};
    \item \texttt{TYPE\_NAME=COMMAND};
    \item \texttt{MODE\_CODE=1};
    \item \texttt{RT=3}.
\end{itemize}

El sistema funciona tambien con el filtro vacio, mostrando toda la informacion disponible.

\section{Pruebas realizadas y observaciones}

\subsection{Prueba integral de funcionamiento}

Se realizo una prueba integral conectando:

\begin{itemize}
    \item \textbf{MIL-M} como emisor;
    \item \textbf{MIL-E} como receptor;
    \item \textbf{Flask-mil} como visualizador sobre el puerto serial del esclavo.
\end{itemize}

El resultado observado fue positivo: el dashboard comenzo a mostrar trafico, indicando que el flujo
master -> SPI -> slave -> serial -> Flask estaba funcional.

\subsection{Warning observado: SEQ\_GAP}

Durante las pruebas se observo al menos un warning \texttt{SEQ\_GAP}. Este warning significa que
el receptor esperaba una secuencia consecutiva de mini-tramas, pero detecto un salto en el valor
de \texttt{SEQ}.

En esta etapa del desarrollo, un warning aislado de este tipo al inicio de la captura no se considera
critico. Las causas mas probables son:

\begin{itemize}
    \item enganche del receptor o del Flask a mitad de una transmision ya iniciada;
    \item reinicio de una de las placas;
    \item desfasaje inicial entre master y slave.
\end{itemize}

La interpretacion actual es que un \texttt{SEQ\_GAP} unico al principio no invalida la prueba si el
resto del flujo permanece estable.

\section{Validaciones implementadas}

Las validaciones implementadas hasta ahora son tanto de transporte como de protocolo logico.

\subsection{Validaciones de transporte}

\begin{longtable}{|>{\raggedright\arraybackslash}p{4.2cm}|>{\raggedright\arraybackslash}p{8.8cm}|}
\hline
\textbf{Validacion} & \textbf{Descripcion} \\
\hline
\texttt{SYNC} & Verifica que el primer byte de la mini-trama sea \texttt{0xAA}. \\
\hline
Paridad & Verifica la paridad impar sobre los 16 bits de palabra. \\
\hline
\texttt{CHECKSUM} & Verifica el XOR de los 6 primeros bytes del frame SPI. \\
\hline
\texttt{SEQ} & Verifica continuidad de secuencia entre palabras consecutivas. \\
\hline
\end{longtable}

\subsection{Validaciones de mensaje}

\begin{longtable}{|>{\raggedright\arraybackslash}p{4.2cm}|>{\raggedright\arraybackslash}p{8.8cm}|}
\hline
\textbf{Chequeo} & \textbf{Descripcion} \\
\hline
Orden de palabras & Verifica si un \texttt{STATUS} o un \texttt{DATA} aparecen en un orden inconsistente respecto del mensaje activo. \\
\hline
Cantidad de datos & Verifica si las palabras de datos vistas son menos o mas que las esperadas. \\
\hline
Broadcast & Verifica si aparece una palabra de estado donde logica y conceptualmente no deberia haber respuesta. \\
\hline
Mode Code & Permite distinguir mensajes especiales de control. \\
\hline
Estado huerfano & Detecta \texttt{STATUS} o \texttt{DATA} fuera de contexto de mensaje. \\
\hline
\end{longtable}

\section{Estado tecnico alcanzado}

En el estado actual del proyecto se puede afirmar que ya se dispone de una simulacion logica
funcional y relativamente fiel de MIL-STD-1553 sobre SPI. Lo alcanzado incluye:

\begin{itemize}
    \item mini-tramas estables de transporte;
    \item palabras \texttt{COMMAND}, \texttt{DATA} y \texttt{STATUS};
    \item mensajes \texttt{BC->RT}, \texttt{RT->BC}, \texttt{BROADCAST} y \texttt{MODE\_CODE};
    \item \texttt{Word Count} variable;
    \item soporte de \texttt{Word Count = 0 => 32} a nivel logico;
    \item deteccion de \texttt{Mode Code};
    \item deteccion de \texttt{Broadcast};
    \item validaciones de transporte;
    \item validaciones de contexto de mensaje;
    \item visualizacion provisional en Flask.
\end{itemize}

\section{Limitaciones actuales}

Persisten aun varias limitaciones, esperables para esta etapa:

\begin{itemize}
    \item la capa fisica sigue siendo SPI y no un bus 1553 real;
    \item el receptor interpreta mensajes logicos simples, pero todavia no maneja \texttt{RT->RT}
    mediado por BC;
    \item no se implemento aun un paquete host mas rico por palabra;
    \item el dashboard actual es una herramienta provisional y no el host final;
    \item el contador \texttt{SEQ} utiliza un solo byte;
    \item no existe aun timestamp fino por palabra al estilo del informe tecnico 2.
\end{itemize}

\section{Respuesta a una pregunta clave del proyecto}

Una de las preguntas centrales durante esta etapa fue si todo este trabajo ya puede considerarse
una simulacion de mini-tramas tipo STD-MIL con validacion. La respuesta actual es:

\begin{quote}
Si, a nivel logico y de transporte, el sistema ya puede considerarse una simulacion funcional de
mini-tramas inspiradas en MIL-STD-1553 con su respectiva validacion.
\end{quote}

Esto no significa aun compatibilidad fisica con un bus real, pero si representa una base valida y
coherente para seguir evolucionando sin cambiar toda la arquitectura.

\section{Proximos pasos recomendados}

Los siguientes pasos tecnicamente razonables son:

\begin{enumerate}
    \item incorporar soporte logico para mensajes \texttt{RT->RT} mediados por BC;
    \item adaptar el receptor para interpretar correctamente esa secuencia;
    \item enriquecer el dashboard Flask con filtros mas especificos por RT, SA y Mode Code;
    \item definir el formato final de paquete host por palabra para la Raspberry dedicada;
    \item migrar el host provisional a la Raspberry final cuando el hardware este disponible;
    \item acercar gradualmente el sistema a la capa fisica real 1553.
\end{enumerate}

\section{Conclusion}

El trabajo realizado hasta el momento permitio transformar un enlace SPI simple en una plataforma
de simulacion logica de MIL-STD-1553 mucho mas ordenada, trazable y util para continuar el
desarrollo. Se consolidaron componentes separados, un protocolo comun, mensajes mas realistas,
validaciones de integridad y de contexto, y una visualizacion temporal en Flask que ya permite
comprobar experimentalmente el funcionamiento de la cadena completa.

En terminos practicos, el proyecto se encuentra en un punto muy valioso: todavia es simple de
depurar, pero ya representa suficiente semantica del estandar como para que el siguiente salto,
hacia una implementacion fisica mas realista, pueda hacerse con mucho menor riesgo.

\end{document}
